\par\noindent\textbf{(1)} 设 \(A_2(0)=0\)。沿长度 \(L\) 积分得 \[A_2(L)=i\kappa A_1^2\int_0^L e^{i\Delta kz}dz=i\kappa A_1^2 L e^{i\Delta kL/2}\operatorname{sinc}(\Delta kL/2)\]，其中 \(\operatorname{sinc}x=\sin x/x\)。
\par\noindent\textbf{(2)} 倍频功率与 \(|A_2|^2\) 成正比，故 \(P_{2\omega}\propto|\kappa|^2|A_1|^4L^2\operatorname{sinc}^2(\Delta kL/2)\)。完美相位匹配 \(\Delta k=0\) 时随 \(L^2\) 增长，而失配时出现振荡零点，反映不同晶体位置生成波的相干叠加。
\par\noindent\textbf{(3)} 源项相位相差 \(\pi\) 的长度为相干长度 \(L_c=\pi/|\Delta k|\)。当晶长超过这一尺度，后段生成的倍频场可能抵消前段；\(\Delta kL=2\pi n\)（非零整数）时未耗尽近似下倍频幅恰为零。
\par\noindent\textbf{(4)} 在准相位匹配中按空间周期 \(\Lambda\) 翻转非线性系数，利用倒格矢 \(G=2\pi/\Lambda\) 补偿失配，使 \(\Delta k-G\approx0\)，但有效非线性系数由傅里叶谐波降低。强转换时 \(A_1(z)\) 不再常数，应联立Manley–Rowe能量方程，不能保持无限二次增长。
